\section{Research extensions}

The correction itself is expository. A genuine research program can begin by extending the diagnostic framework beyond canonical flat phase space.

\subsection{Noncanonical Poisson manifolds}

For general local coordinates \(z^a\), let
\begin{equation}
\pb{z^a}{z^b}=\Pi^{ab}(z),
\end{equation}
where \(\Pi^{ab}\) is an antisymmetric Poisson tensor satisfying the Jacobi identity. A deformation quantization has the local form
\begin{equation}
f\star g
=fg+\frac{\ii\hbar}{2}\Pi^{ab}(z)
\partial_a f\,\partial_b g+O(\hbar^2),
\end{equation}
and hence
\begin{equation}
\starcomm{f}{g}
=
\ii\hbar\pb{f}{g}+O(\hbar^2).
\label{eq:generalstar}
\end{equation}
Unlike the canonical Moyal product, higher-order terms can depend on derivatives of \(\Pi^{ab}(z)\) and on the chosen equivalent star-product representative.

\subsection{Magnetic, constrained, and curved phase spaces}

Several extensions are mathematically substantive:
\begin{itemize}
  \item magnetic symplectic forms, in which field-strength terms modify the canonical two-form and momentum brackets;
  \item constrained Hamiltonian systems, where Dirac brackets replace canonical Poisson brackets;
  \item Berry-curvature corrections in semiclassical dynamics;
  \item cotangent bundles over curved configuration manifolds;
  \item noncommutative coordinate algebras and deformation-quantized Poisson manifolds.
\end{itemize}
These settings make clear why no fixed cross-gradient rule can replace the underlying geometric data.

\subsection{Symbolic consistency checking}

A practical computational contribution would be an automated checker that accepts a proposed identity and reports:
\begin{enumerate}[label=(\alph*)]
  \item tensor-rank and scalar/vector mismatches;
  \item incompatible physical dimensions;
  \item undefined cross products outside three dimensions;
  \item confusion between commutators, wedge products, and curls;
  \item missing operator-ordering prescriptions;
  \item omitted singular or gauge terms.
\end{enumerate}
Such a tool would convert the present conceptual taxonomy into a reproducible method for screening speculative phase-space formulae.
